\newpage
% ============================================================
%  SECTION: Technologies Used
% ============================================================

\section{Technologies}
\label{sec:technologies}

This section provides a brief overview of each technology that constitutes the
experimental infrastructure used in this thesis. Together these components form
a reproducible, cloud-native testbed for evaluating the two autoscaling
strategies under investigation.

% ------------------------------------------------------------
\subsection{K3s}
\label{subsec:k3s}

K3s is a CNCF-certified, lightweight Kubernetes distribution developed by
Rancher. It packages the full Kubernetes control plane into a single binary of
less than 70~MB and eliminates several heavyweight dependencies of upstream
Kubernetes, making it well-suited for resource-constrained environments such as
edge devices, IoT deployments, and developer workstations (\cite{k3s:docs}).
Despite its reduced footprint, K3s is fully conformant with the Kubernetes API
and supports the complete set of Kubernetes primitives required by this thesis,
including \texttt{StatefulSets}, \texttt{CustomResourceDefinitions}, and Helm
chart deployments. K3s serves as the cluster management layer for all
infrastructure described in this thesis.

% ------------------------------------------------------------
\subsection{FastAPI}
\label{subsec:fastapi}

FastAPI is a modern, high-performance Python web framework designed for
building RESTful APIs. It is built on top of \textit{Starlette} for its
web-server layer and \textit{Pydantic} for data validation, and it uses Python
type hints to enable automatic request validation and interactive OpenAPI
documentation generation (\cite{fastapi:docs}). FastAPI is built around the
Asynchronous Server Gateway Interface (ASGI), enabling non-blocking I/O
operations that are especially advantageous in workloads involving concurrent
database or cache interactions (\cite{fastapi:async}). In the context of this
thesis, a FastAPI application serves as the \ac{HTTP} entrypoint to the cluster,
receiving incoming requests from the external load generator and routing read
and write operations to the Redis Cluster backend.

% ------------------------------------------------------------
\subsection{Redis Cluster}
\label{subsec:redis}

Redis is an open-source, in-memory data structure store widely used as a
cache, session store, and real-time data platform (\cite{redis:io}). The Redis
Cluster topology used in this thesis distributes data across multiple primary
nodes using a hash-slot-based sharding mechanism. The full keyspace is divided
into \num{16384} hash slots, each assigned to a primary node, and each primary
is optionally backed by one or more replica nodes for fault tolerance
(\cite{redis:clusterspec}). Redis Cluster provides high availability through
automatic failover: if a primary becomes unreachable, one of its replicas is
promoted. In this thesis, the Redis Cluster is deployed as a Kubernetes
\texttt{StatefulSet} with a minimum of three primary nodes and forms the core
stateful workload under performance evaluation.

% ------------------------------------------------------------
\subsection{Horizontal Pod Autoscaler (HPA)}
\label{subsec:hpa}

The Horizontal Pod Autoscaler is a built-in Kubernetes controller that
automatically adjusts the number of pod replicas in a workload based on
observed resource metrics, most commonly CPU and memory utilisation
(\cite{kubernetes:autoscaling}). The HPA periodically queries the Kubernetes
Metrics Server (or a custom metrics adapter) and computes a desired replica
count proportional to the ratio of current to target utilisation. Its
synchronisation loop runs every 15~seconds by default. While effective for
many stateless use cases, the HPA is a fundamentally \textit{reactive}
mechanism: scaling decisions are always downstream of observed resource
conditions, making it vulnerable to the lag issues described in
Section~\ref{subsec:stateful-challenges}. In this thesis, the HPA is deployed
as the \textit{naive} baseline scaling strategy, configured to scale the Redis
\texttt{StatefulSet} in response to CPU and memory utilisation thresholds.

% ------------------------------------------------------------
\subsection{KEDA}
\label{subsec:keda}

Kubernetes Event-Driven Autoscaling (KEDA) is a CNCF-graduated open-source
project that extends the Kubernetes HPA with the ability to scale workloads
based on event sources and external metrics beyond simple CPU and memory
(\cite{keda:home}). KEDA introduces the \texttt{ScaledObject} custom resource,
which declares a scaling target (such as a \texttt{Deployment} or
\texttt{StatefulSet}), a polling interval, cooldown periods, and one or more
\textit{scalers} that define the metric source (\cite{keda:scaledobjectspec}).
When the configured metric from the scaler exceeds a threshold, KEDA adjusts
the replica count of the target workload by updating the underlying HPA on the
user's behalf. In this thesis, KEDA is configured with a Prometheus scaler
that reads application-level metrics — such as request queue depth and cache
hit ratio — directly from the Prometheus monitoring stack, enabling proactive,
context-aware scaling decisions before resource exhaustion occurs
(\cite{keda:prometheusscaler}).

% ------------------------------------------------------------
\subsection{Prometheus and Grafana}
\label{subsec:observability}

The observability stack in this thesis is deployed using the
\texttt{kube-prometheus-stack} Helm chart, which bundles Prometheus, Grafana,
and Alertmanager into a single, opinionated release (\cite{artifacthub:kube-prometheus-stack}).
\textit{Prometheus} is an open-source monitoring system and time-series
database that collects metrics by periodically scraping HTTP endpoints exposed
by instrumented targets (\cite{prometheus:overview}). In the Kubernetes context,
it discovers scrape targets automatically via Kubernetes service discovery,
collecting cluster-level metrics from \texttt{kube-state-metrics} and
\texttt{node-exporter}, as well as application-level metrics from the FastAPI
and Redis exporters. \textit{Grafana} is an open-source analytics and
visualisation platform that queries Prometheus via its native data source
integration and renders the collected time-series data as configurable
dashboards (\cite{grafana:github}). Grafana dashboards are used throughout the
experimental evaluation in this thesis to visualise autoscaling behaviour,
latency distributions, and resource utilisation over time. Crucially, the
Prometheus instance also acts as the external metric backend for the KEDA
Prometheus scaler, creating a unified data path from application telemetry to
autoscaling decisions.

% ------------------------------------------------------------
\subsection{Helm}
\label{subsec:helm}

Helm is the de-facto package manager for Kubernetes. It introduces the
concept of a \textit{chart} — a collection of templated Kubernetes manifests
that describe a deployable application — and a \textit{release}, which is a
running instance of a chart in a cluster (\cite{helm:docs}). Helm charts enable
repeatable, version-controlled deployments with configurable values overrides
and built-in rollback support, significantly reducing the operational
complexity of deploying multi-resource applications such as the
\texttt{kube-prometheus-stack}. In this thesis, Helm is used to deploy the
observability stack, and the deployment configurations are parameterised
through values files that are committed alongside the thesis infrastructure
code.

% ------------------------------------------------------------
\subsection{Grafana k6}
\label{subsec:k6}

Grafana k6 is an open-source, developer-focused performance and load testing
tool (\cite{k6:docs}). Tests in k6 are written as JavaScript modules, and k6
executes them using a Go runtime which maps each \textit{virtual user} (VU)
to a lightweight goroutine, enabling high concurrency with low resource
overhead (\cite{k6:github}). k6 supports configurable test scenarios including
ramp-up and ramp-down of virtual users over time, making it possible to
simulate realistic traffic patterns such as gradual load increases, sustained
peaks, and sudden spikes (\cite{k6:apitesting}). In the context of this thesis,
k6 runs \textit{outside} the Kubernetes cluster as an independent load
generator. It sends configurable volumes of HTTP traffic through the FastAPI
entrypoint and into the Redis Cluster, producing the synthetic workload
against which both autoscaling strategies — HPA and KEDA — are evaluated.
The ability to precisely control the number of virtual users and the ramp
profile makes k6 well-suited for generating reproducible, parameterised load
scenarios that can isolate the performance characteristics of each scaling
approach.
